\scanpage{52}
% manuscript line(s) 1–3
\[
\deleted{\unsure{N(b|\mathrm{N.B.}}}\quad\begin{aligned}&\left\{\begin{aligned}T(b,\unsure{r|\gamma})&=b\\ \delta(b,r)&=\deleted{\unsure{b}}r^2\cdot\delta\\ N(b,r)&=(T^2\deleted{\unsure{+}}\unsure{-|+}\delta)(b,r)/4=(b^2\deleted{\unsure{+}}-r^2\cdot\delta)/4\end{aligned}\right. \end{aligned}
\]
% manuscript line(s) 3
\margin{donc}
% manuscript line(s) 4
C'est une section de \deleted{\unsure{$\Gamma$|de}}$\underline O_X$, car
% manuscript line(s) 5
\[
b\equiv rT\ (2)\qquad\text{donc}\qquad b^2\equiv r^2T^2\ (4)
\]
% manuscript line(s) 5–7
% slanted
% manuscript line(s) 5–7
\margin{\uline{NB} On en déduit $\varphi_N((b,r),(b',r'))=\dfrac{bb'\deleted{\unsure{+}}\inserted{-}rr'\cdot\delta}{\deleted{\unsure{4}}2}$}
% manuscript line(s) 6
donc
% manuscript line(s) 7
\[
b^2\deleted{\unsure{+}}-r^2\cdot\delta\unsure{\simeq|\equiv}r^2(\underbrace{T^2\deleted{\unsure{+}}-\delta}_{\equiv0\ (4)})\equiv0\ (4)
\]

% manuscript line(s) 8
Soit $\underline{x\in(b,r)}\in\Gamma\mathcal A$, on a
% manuscript line(s) 9
\[
x^2\deleted{\unsure{+}}-bx\deleted{\unsure{-}}+\frac{b^2\deleted{\unsure{+}}-r^2\cdot\delta}{4}=0\qquad\text{i.e.}
\]
% manuscript line(s) 10–11
\[
\begin{aligned}x^2&=bx-\unsure{\cdot}\frac{b^2-r^2\cdot\delta}{4}=\left(b^2-\frac{b^2-r^2\delta}{2},br\right)\\ &\phantom{=bx-\frac{b^2-r^2\cdot\delta}{4}}=\left(\frac{b^2+r^2\cdot\delta}{2},br\right)\end{aligned}
\]
% manuscript line(s) 12
Cela permet de calculer, si
% manuscript line(s) 13
\[
x=(b,r),\qquad x'=(b',r')
\]
% manuscript line(s) 14
le produit $xx'$ comme
% manuscript line(s) 15
\[
xx'=\tfrac12\big((x+x')^2-x^2-x'^2\big)
\]
% manuscript line(s) 16–18
\[
\begin{aligned}\text{i.e.}\quad 2xx'=\Big(&\frac{(b+b')^2+(r\unsure{-|+}r')^2\cdot\delta}{2}-\deleted{\unsure{b^2}}\frac{b^2+r^2\cdot\delta}{2}-\frac{b'^2+r'^2\cdot\delta}{2},\\ &(b+b')(r+\deleted{\unsure{b'}}r')\deleted{\unsure{r}}-br-b'r'\Big)\Big)\\ =\Big(&\underbrace{bb'+rr'\cdot\delta}_{\substack{\mapsto r_0r'_0\delta_0+r_0r'_0\delta_0\\=2r_0r'_0\delta_0=0}},\underset{\substack{r_0r'_0T_0+r_0r'_0T_0\\=2r_0r'_0T_0=0}}{br'+b'r}\Big)\end{aligned}
\]
% manuscript line(s) 18–20
\margin{mod $2$ c'est bien nul car $b_0=r_0T_0$, $b'_0=r'_0T_0$ $T_0^2=\delta_0$}
% manuscript line(s) 20
donc
% manuscript line(s) 21
\[
\boxed{xx'=(b,r)(b',r')=\left(\tfrac12(bb'\deleted{\unsure{-}}+rr'\cdot\delta,\tfrac12(br'+rb')\right)\big)}
\]
% archive sheet number: 51
